\documentclass{article}

\usepackage{amsmath, makecell}

\title{Ch. 25 HW \#2 Solutions}
\author{Principles of Macroeconomics}
\date{\today}

\begin{document}

\maketitle

\section*{Problems and Applications}

\begin{enumerate}

\setcounter{enumi}{6}

\item

	\begin{enumerate}
	
	\item $\frac{1.47 - 0.88}{0.88} \cdot 100 = 67\%$
	
	\item $\frac{25.86 - 6.57}{6.57} \cdot 100 = 294\%$
	
	\item 
		\begin{gather*}
		1980: \$0.88 \cdot \frac{1 \text{ hour}}{\$6.57} \cdot \frac{60 \text{ minutes}}{1 \text{ hour}} = 8.04 \text{ minutes} \\
		2021: \$1.47 \cdot \frac{1 \text{ hour}}{\$25.86} \cdot \frac{60 \text{ minutes}}{1 \text{ hour}} = 3.41 \text{ minutes}
		\end{gather*}
		
	\item It rose because they could buy more after an hour of labor.
	
	\end{enumerate}
	
\item

	\begin{enumerate}
	
	\item It does. Since Social Security benefits increase in proportion to the CPI and the CPI overstates inflation, Social Security benefits increase faster than the price level. That means senior citizens' purchasing power increases.

	\item I would determine a fixed basket that more accurately represents senior citizens' consumption. Then, I would calculate a new measure of the price level using the senior citizen specific fixed basket. Lastly, I would see if the new measure of the price level was increasing faster than the CPI. If not, then senior citizens would be better off from year to year. If so, then they would be worse off from year to year. 
	
	\end{enumerate}
	
\item

	\begin{enumerate}
	
	\item The real interest rate is the nominal interest rate minus the inflation rate. If the inflation rate is higher than expected, then the real interest rate is lower than expected.
	
	\item The lender loses because their purchasing power is increasing by less than expected. The borrower gains because they're giving up less purchasing power than expected.
	
	\item If inflation was unexpectedly high and the nominal interest rate was fixed, then the real interest rate was unexpectedly low. That benefits the homeowners (the borrowers) and hurts the banks (the lenders) for the reasons explained in (b).
	
	\end{enumerate}

\end{enumerate}

\section*{Additional Problems}

\begin{enumerate}

\item  The family gets worse off from 2020-2025.
	\begin{gather*}
	2020 : \$40,000 \cdot \frac{100}{100} = \$40,000 \\
	2025 : \$44,000 \cdot \frac{100}{125} = \$35,200 
	\end{gather*}
	
\item

	\begin{enumerate}
	
	\item \begin{equation*}
		\$14.35 \left( \frac{100}{292.7} \right) = \$4.90
		\end{equation*}
		
	\item \begin{gather*}
		\$4.90 = 0.95 \cdot (\text{Real entry-level wage in 2015}) \\
		\frac{\$4.90}{0.95} = \text{Real entry-level wage in 2015} \\
		\$5.16 = \text{Real entry-level wage in 2015}
		\end{gather*}
		
	\item \begin{equation*}
		\$5.16 \left( \frac{237}{100} \right) = \$12.23
		\end{equation*}
		
	\item \begin{equation*}
		\left( \left( \frac{292.7}{237} \right)^{\frac{1}{7}} - 1 \right) \cdot 100 = 3.06\%
		\end{equation*}
	
	\end{enumerate}
	
\end{enumerate}

\end{document}